% page 456 du fac-simile (image p-455 du scan)
% bloc 147.3 x 250.6 mm ; marge gauche 8.0 mm
\pgc{-12.700mm}{0.000mm}{
\l{\cel{3}448}
\l{}
\l{\sou{plumo}. -I. (\sou{zool.)}\cel{1}Singla ek la tubi korno-harda, garni-}
\l{\cel{3}sita ye aristi e ye lanugo, qua kovras la korpo di la}
\l{\cel{3}uceli. - II. Tubo di plumo, taliita por skribar perlu}
\l{\cel{3}(pos trempir lua pinto aden inko). - III. Mikra lamo}
\l{\cel{3}de metalo, qua havas la formo di la pinto di plumo}
\l{\cel{3}taliita, di qua onu glite muntas la extremajo sen-pinta}
\l{\cel{3}an-en tubo ucel-plumatra. - EFIS.}
\l{}
\l{\sou{plump}\cel{1}!\cel{2}Interjeciono qua expresas la bruiso de falo, pre-}
\l{\cel{3}cipue di ulo tre pezoza, mola, inerta.}
\l{}
\l{\sou{plumpa}. (\sou{pri vivanto}). Qua su movas quale se lu esus tre}
\l{\cel{3}pezoza, pro la denseso di sua korpo. - D.}
\l{}
\l{\sou{plunjar}\cel{1}(\sou{netrans}.) Penetrar, sinkar profunde aden aquo,}
\l{\cel{3}propra-vole. -\cel{2}EF.}
\l{}
\l{\sou{plura}. Ula quanto (ne tre granda) de personi, de kozi -}
\l{\cel{3}mem nur du. - DEFIS.}
\l{}
\l{\sou{plus}. I. Ultre to ; adjuntenda a... - II. (\sou{kun \textquotedbl{}ne\textquotedbl{} avan}}
\l{\cel{3}od\cel{1}\sou{ante}\cel{1}lu) - indikas la ceso di la ago, di la eso.}
\l{}
\l{\sou{plusho}. Stofo simila a veluro, e kun longa pili e lanugo}
\l{\cel{3}sur la averso. - DEFR.}
\l{}
\l{\sou{plusquamperfekto}. (\sou{gram.)}\cel{2}Verbo-tempo qua indikas pasinto}
\l{\cel{3}antea ye altra tempo pasinta (t. e. ...is...inta). -}
\l{}
\l{}
\l{\sou{plutokratio}. Guverno da la richi. - DEFIRS.}
\l{}
\l{\sou{plutonala}. (\sou{geol.}) (\sou{pri rok}i) Produktita da la efiko di}
\l{\cel{3}fairo sub-ter-sula naturala. - DEFIRS.}
\l{}
\l{\sou{pluvar}.\cel{2}(\sou{netr.}) (\sou{aludante la aquo-vaporo di la nubi}).}
\l{\cel{3}Kondensesar e falar gutope.FIS.}
\l{}
\l{\sou{pluviometro.}\cel{1}(\sou{fiziko)}. Instrumento quan onu uzas por me-}
\l{\cel{3}zurar la quanto de pulv-aquo qua falas en un loko dum}
\l{\cel{3}ula quanto determinita de tempo. - DEFIS.}
\l{}
\l{\sou{pneumatika.}\cel{1}Dicesas pri mashino o pumpo per qua onu pumpas}
\l{\cel{3}la aero di tanko por eventigar lua vakueso. - DEFIS.}
\l{}
\l{\sou{pneumokoko.}\cel{1}(\sou{patol.})Agento patogena, la maxim ordinara, di}
\l{\cel{3}pulmonito, deskovrita da Louis Pasteur en la yaro 1881.}
\l{}
\l{\cel{1}po.\cel{2}Kom preco di... - Kambie di...}
\l{}
\l{\cel{1}\sou{pociono.}\cel{1}Medikamento quan onu glutas forme di drinkajo.}
\l{\cel{3}EF.}
\l{}
\l{\cel{1}\sou{podagro}. (\sou{patol.}) Morbo qua afektas, ordinar-kaze, la ar-}
\l{\cel{4}tiki, iras de ica ad ita, e quan onu imputis olim a la}
\l{\cel{4}influo da guti de humoro. -}
}
